% =====================================================================
% PART 5: ECONOMICS & SUSTAINABILITY
% =====================================================================
\section{Economics \& Sustainability}

% --- 5.1 Key Question ---
\begin{frame}{Key Question}
\begin{center}
\Large\bfseries
How much does it \emph{really} cost to train and serve an LLM---\\[4pt]
in dollars, kilowatt-hours, and tonnes of CO$_2$?
\end{center}
\vfill
\begin{quotation}
\small\itshape
``A 1024-GPU H100 cluster costs \$30M in hardware. Electricity is
surprisingly small ($\sim$10\% of total). The dominant cost lever is
utilization---idle GPUs cost the same as busy ones.''
\end{quotation}
\end{frame}

% --- 5.2 The TCO Equation ---
\begin{frame}{Wall 17: The Capital Wall (TCO)}
\begin{block}{The Equation}
\[
\text{TCO} = \underbrace{\text{CapEx}}_{\text{hardware + facility}}
          + \underbrace{\text{OpEx}_{\text{energy}}}_{\text{electricity}}
          + \underbrace{\text{OpEx}_{\text{maint}}}_{\text{staff, repairs}}
\]
\end{block}
\vspace{0.5em}
\textbf{Key insight:} GPUs are only $\sim$40\% of total CapEx.\\
Networking, cooling, facility, and staff add 50--150\%.

\vspace{0.5em}
\begin{itemize}
  \item \texttt{infrastructure\_multiplier = 2.0--2.5} for full datacenter TCO
  \item Amortized over 3--5 year depreciation schedule
  \item Cloud spot instances can beat on-prem for bursty workloads
\end{itemize}
\end{frame}

% --- 5.3 Infrastructure Multiplier ---
\begin{frame}{The Infrastructure Multiplier}
\centering
\begin{tabular}{lrr}
\toprule
\textbf{Component} & \textbf{Cost (\%)} & \textbf{Multiplier} \\
\midrule
GPU Hardware       & 40\% & 1.0$\times$ \\
Network (IB)       & 20\% & --- \\
Power \& Cooling   & 15\% & --- \\
Facility / Land    & 10\% & --- \\
Staff / Operations & 15\% & --- \\
\midrule
\textbf{Total}     & 100\% & \textbf{2.5$\times$} \\
\bottomrule
\end{tabular}

\vspace{1em}
\alert{Pitfall:} Budgeting only for GPUs underestimates TCO by 2--3$\times$.
\end{frame}

% --- 5.4 The Sustainability Equation ---
\begin{frame}{Wall 18: The Sustainability Wall}
\begin{block}{The Equation}
\[
\text{Carbon (kg CO}_2\text{)} = \text{Energy (kWh)} \times \text{PUE} \times \text{Carbon Intensity (g/kWh)}
\]
\end{block}
\vspace{0.5em}

\textbf{Three levers:}
\begin{enumerate}
  \item \textbf{Energy} --- Reduce compute time (better MFU, smaller model)
  \item \textbf{PUE} --- Better cooling (air 1.3 $\to$ liquid 1.06)
  \item \textbf{Carbon intensity} --- Choose your grid (Quebec 20 vs Poland 820 g/kWh)
\end{enumerate}

\vspace{0.5em}
\begin{center}
\alert{Geography is the single biggest lever for sustainable AI.}
\end{center}
\end{frame}

% --- 5.5 Geography Matters ---
\begin{frame}{Carbon Intensity: The 41$\times$ Gap}
\centering
\begin{tabular}{lcr}
\toprule
\textbf{Region} & \textbf{Grid Mix} & \textbf{gCO$_2$/kWh} \\
\midrule
Quebec, Canada    & Hydro          & 20 \\
Sweden            & Hydro + Nuclear & 25 \\
US Average        & Mixed          & 390 \\
Germany           & Coal + Wind    & 380 \\
Poland            & Coal           & 820 \\
\bottomrule
\end{tabular}

\vspace{1em}
Training the \emph{same model} on the \emph{same hardware}:\\[4pt]
\textbf{Poland:} $\sim$800 tonnes CO$_2$ \quad vs \quad \textbf{Quebec:} $\sim$20 tonnes CO$_2$

\vspace{0.5em}
$\Rightarrow$ \textbf{41$\times$} difference---more impactful than any algorithmic optimization.
\end{frame}

% --- 5.6 Embodied Carbon ---
\begin{frame}{Embodied Carbon: Manufacturing Dominates at the Edge}
\begin{columns}[T]
\column{0.5\textwidth}
\textbf{Cloud / Training}
\begin{itemize}
  \item Operational carbon dominates
  \item GPU runs 24/7 for months
  \item Embodied: $\sim$5\% of lifecycle
\end{itemize}

\column{0.5\textwidth}
\textbf{IoT / TinyML}
\begin{itemize}
  \item Manufacturing carbon dominates
  \item MCU uses microwatts
  \item Embodied: $>$90\% of lifecycle
  \item Multiplied by millions of devices
\end{itemize}
\end{columns}

\vspace{1em}
\begin{center}
\small
Source: Gupta et al.\ (2022), ``ACT: Designing Sustainable Computer\\
Systems with an Architectural Carbon Modeling Tool''
\end{center}
\end{frame}

% --- 5.7 Live Demo: Economics ---
\begin{frame}[fragile]{Live Demo: TCO + Carbon Analysis}
\begin{minted}[fontsize=\scriptsize]{python}
import mlsysim

fleet = mlsysim.Systems.Clusters.Research_256  # 256x H100
econ  = mlsysim.EconomicsModel()

# Full TCO with 2.5x infrastructure multiplier
result = econ.solve(
    fleet=fleet,
    duration_days=30,
    datacenter=mlsysim.Infra.Grids.Quebec,
    mfu=0.45,
    infrastructure_multiplier=2.5,
    amortization_years=3.0
)

print(f"CapEx (30-day amortized): ${result.capex_usd:,.0f}")
print(f"OpEx (energy):            ${result.opex_energy_usd:,.0f}")
print(f"OpEx (maintenance):       ${result.opex_maintenance_usd:,.0f}")
print(f"Total TCO:                ${result.tco_usd:,.0f}")
print(f"Carbon:                   {result.carbon_footprint_kg/1000:.1f} tonnes")
\end{minted}
\end{frame}

% --- 5.8 Live Demo: Geography Comparison ---
\begin{frame}[fragile]{Live Demo: The Carbon Geography Experiment}
\begin{minted}[fontsize=\scriptsize]{python}
import mlsysim

fleet = mlsysim.Systems.Clusters.Research_256
solver = mlsysim.SustainabilityModel()

for region in [mlsysim.Infra.Grids.Poland, mlsysim.Infra.Grids.Quebec]:
    result = solver.solve(fleet=fleet, duration_days=30,
                          datacenter=region, mfu=0.45)
    print(f"{result.region_name:>12}: "
          f"{result.total_energy_kwh.to('MWh'):.0f} MWh, "
          f"{result.carbon_footprint_kg/1000:.1f} tonnes CO2")
\end{minted}

\vspace{0.5em}
\begin{exampleblock}{Expected Output}
\ttfamily\scriptsize
\hspace{1em}Poland: 1,064 MWh, 872.5 tonnes CO2\\
\hspace{1em}Quebec: 1,064 MWh, \phantom{0}21.3 tonnes CO2
\end{exampleblock}

\vspace{0.3em}
\alert{Same energy. 41$\times$ less carbon. Geography wins.}
\end{frame}

% --- 5.9 Exercise ---
\begin{frame}[fragile]{Exercise: Carbon-Aware Placement}
\begin{alertblock}{Your Task (5 minutes)}
Using the \texttt{SustainabilityModel}, compare a 30-day training run
on 512 H100s in \textbf{Poland} vs \textbf{Quebec}. How many tonnes of
CO$_2$ do you save by choosing Quebec?
\end{alertblock}

\begin{minted}[fontsize=\scriptsize]{python}
import mlsysim
from mlsysim.infra.types import Datacenter

fleet = mlsysim.Fleet(
    name="Training Fleet", node=mlsysim.Systems.Nodes.DGX_H100,
    count=64, fabric=mlsysim.Systems.Fabrics.InfiniBand_NDR)
solver = mlsysim.SustainabilityModel()
# YOUR CODE: solve for Poland and Quebec, compute the difference
\end{minted}

\note{Expected answer: $\sim$850 tonnes CO$_2$ saved.
The energy is identical ($\sim$2{,}100 MWh)---only carbon intensity differs.
Key insight: the cheapest carbon reduction is a \texttt{git push} to a different region.}
\end{frame}

% --- 5.10 Takeaway ---
\begin{frame}{Key Takeaway: Economics \& Sustainability}
\begin{center}
\Large
\begin{enumerate}
  \item \textbf{TCO $\neq$ GPU cost.} Infrastructure is 2--2.5$\times$.
  \item \textbf{Utilization is the cost lever.} Idle GPUs cost the same.
  \item \textbf{Geography $>$ algorithms} for carbon reduction.
\end{enumerate}

\vspace{1em}
\normalsize
\textit{``The cheapest watt is the one you don't consume.\\
The cleanest watt is the one from a hydro dam.''}
\end{center}
\end{frame}


% =====================================================================
% PART 6: DESIGN SPACE EXPLORATION
% =====================================================================
\section{Design Space Exploration}

% --- 6.1 Key Question ---
\begin{frame}{Key Question}
\begin{center}
\Large\bfseries
Given a budget and an SLA,\\[4pt]
how do you find the \emph{best} hardware configuration?
\end{center}
\vfill
\small
The search space is combinatorial:
$|\text{hardware}| \times |\text{batch sizes}| \times |\text{precisions}| \times |\text{parallelism configs}|$\\
can easily exceed $10^4$ configurations.
\end{frame}

% --- 6.2 The DSE Pattern ---
\begin{frame}{The DSE Pattern: Declare, Search, Rank}
\begin{enumerate}
  \item \textbf{Declare} the search space:
    \begin{itemize}
      \item Hardware: \{A100, H100, H200, B200\}
      \item Batch sizes: \{1, 2, 4, 8, 16, 32, 64\}
      \item Precisions: \{fp16, int8, int4\}
    \end{itemize}
  \item \textbf{Search} with an objective:
    \begin{itemize}
      \item \texttt{minimize: tco\_usd}
      \item \texttt{maximize: throughput}
    \end{itemize}
  \item \textbf{Rank} candidates subject to constraints:
    \begin{itemize}
      \item \texttt{latency < 50 ms}
      \item \texttt{feasible == True}
    \end{itemize}
\end{enumerate}

\vspace{0.5em}
mlsysim's \texttt{DSE} class implements exhaustive grid search.\\
Analytical models $\Rightarrow$ each evaluation takes $<$1\,ms.
\end{frame}

% --- 6.3 Pareto Fronts ---
\begin{frame}{Pareto Fronts: No Free Lunch}
\centering
\includegraphics[width=0.65\textwidth]{figures/pareto-placeholder.pdf}

\vspace{0.5em}
\textbf{Pareto front}: the set of configurations where\\
improving one metric \emph{must} worsen another.

\begin{itemize}
  \item Lower latency $\Rightarrow$ lower throughput (smaller batch)
  \item Higher throughput $\Rightarrow$ higher cost (more GPUs)
  \item The ``knee'' of the curve is usually the sweet spot
\end{itemize}
\end{frame}

% --- 6.4 Live Demo: Engine.sweep ---
\begin{frame}[fragile]{Live Demo: Design Space Sweep}
\begin{minted}[fontsize=\scriptsize]{python}
import mlsysim

model = mlsysim.Models.Language.Llama3_8B
hardware_list = [mlsysim.Hardware.A100, mlsysim.Hardware.H100,
                 mlsysim.Hardware.H200, mlsysim.Hardware.B200]

# Sweep over hardware x batch_size x precision
results = mlsysim.Engine.sweep(
    model, hardware_list,
    batch_sizes=[1, 4, 16, 64],
    precisions=["fp16", "int8"],
    efficiency=0.5
)

for r in sorted(results, key=lambda x: x['profile'].latency.magnitude):
    p = r['profile']
    print(f"{r['hardware']:>15} bs={r['batch_size']:<3} {r['precision']:<5} "
          f"| {p.latency.to('ms'):>10.2f~P} | {p.bottleneck:<8} "
          f"| MFU={p.mfu:.3f}")
\end{minted}
\end{frame}

% --- 6.5 Live Demo: DSE with Constraints ---
\begin{frame}[fragile]{Live Demo: DSE with Objective \& Constraints}
\begin{minted}[fontsize=\scriptsize]{python}
from mlsysim.core.dse import DSE

dse = DSE(
    space={
        "batch_size": [1, 4, 8, 16, 32, 64],
        "precision": ["fp16", "int8"],
    },
    objective="maximize: throughput",
    constraints=["latency < 100"]  # ms
)

def evaluate(params):
    p = mlsysim.Engine.solve(
        mlsysim.Models.Language.Llama3_8B,
        mlsysim.Hardware.H100,
        batch_size=params["batch_size"],
        precision=params["precision"])
    return {"throughput": p.throughput.magnitude,
            "latency": p.latency.to("ms").magnitude}

best = dse.search(evaluate)
print(f"Best config: {best['best_params']}")
print(f"Throughput:  {best['best_objective']:.1f} seq/s")
\end{minted}
\end{frame}

% --- 6.6 Batching Optimizer ---
\begin{frame}[fragile]{Live Demo: Batching Optimizer (Pareto Front)}
\begin{minted}[fontsize=\scriptsize]{python}
import mlsysim

optimizer = mlsysim.BatchingOptimizer()
result = optimizer.solve(
    model=mlsysim.Models.Language.Llama3_8B,
    hardware=mlsysim.Hardware.H100,
    seq_len=128,
    arrival_rate_qps=10.0,
    sla_latency_ms=20000.0,  # 20s SLA
    precision="fp16"
)

print(f"Optimal batch size: {result.best_batch_size}")
print(f"Max throughput:     {result.max_throughput:.1f} seq/s")

for pt in result.pareto_front:
    print(f"  bs={pt['batch_size']:<4} "
          f"lat={pt['p99_latency'].m_as('ms'):.0f} ms  "
          f"tpt={pt['throughput']:.1f} seq/s")
\end{minted}
\end{frame}

% --- 6.7 Exercise ---
\begin{frame}[fragile]{Exercise: Budget-Constrained Design}
\begin{alertblock}{Your Task (5 minutes)}
Given a \textbf{\$1M budget}, find the best hardware + batch size + precision
configuration for serving Llama-3-8B inference under a 50\,ms latency SLA.
\end{alertblock}

\begin{minted}[fontsize=\scriptsize]{python}
# Hint: Use Engine.sweep across hardware tiers,
# filter by feasibility and latency SLA,
# then pick the highest-throughput configuration
# within budget (check unit_cost in the Hardware registry)
\end{minted}

\note{Expected answer: H100 at INT8, batch size 16--32 gives the best
throughput/dollar under 50\,ms. A100 at FP16 fails the SLA at high batch sizes.
B200 is faster but fewer GPUs fit in budget.}
\end{frame}

% --- 6.8 Takeaway ---
\begin{frame}{Key Takeaway: Design Space Exploration}
\begin{center}
\Large
\begin{enumerate}
  \item \textbf{Analytical models} make exhaustive DSE feasible ($<$1\,s for $10^4$ configs)
  \item \textbf{Pareto fronts} reveal the actual tradeoff structure
  \item \textbf{Constraints first}: filter infeasible, then optimize
\end{enumerate}

\vspace{1em}
\normalsize
\textit{``You cannot optimize what you cannot model.\\
You cannot model what you cannot measure.''}
\end{center}
\end{frame}


% =====================================================================
% PART 7: TINYML TO FRONTIER
% =====================================================================
\section{TinyML to Frontier}

% --- 7.1 Key Question ---
\begin{frame}{Key Question}
\begin{center}
\Large\bfseries
Can the same analytical framework model\\[4pt]
a \$2 microcontroller \emph{and} a \$3M GPU rack?
\end{center}
\vfill
\small
The answer is yes---because the \textbf{Roofline model} is universal.\\
Only the numbers change. The physics is the same.
\end{frame}

% --- 7.2 The Nine Orders of Magnitude ---
\begin{frame}{The 9-Order-of-Magnitude Scale Span}
\centering
\begin{tabular}{lrrrr}
\toprule
\textbf{Device} & \textbf{FLOPS} & \textbf{Memory} & \textbf{BW} & \textbf{TDP} \\
\midrule
nRF52840       & 64 MFLOPS  & 256\,KiB SRAM & 0.26 GB/s  & 15\,mW \\
ESP32-S3       & 500 MFLOPS & 512\,KiB SRAM & 0.96 GB/s  & 400\,mW \\
\rowcolor{gray!15}
H100 SXM       & 989 TFLOPS & 80\,GB HBM3   & 3{,}350 GB/s & 700\,W \\
\bottomrule
\end{tabular}

\vspace{1em}
\begin{columns}[c]
\column{0.3\textwidth}
\centering
\textbf{Compute}\\
$\frac{989 \times 10^{12}}{64 \times 10^{6}} \approx 10^{7}$

\column{0.3\textwidth}
\centering
\textbf{Memory BW}\\
$\frac{3{,}350}{0.000256} \approx 10^{7}$

\column{0.3\textwidth}
\centering
\textbf{Power}\\
$\frac{700}{0.015} \approx 10^{4.7}$
\end{columns}

\vspace{0.5em}
\alert{Same Roofline equation. Nine orders of magnitude apart.}
\end{frame}

% --- 7.3 TinyML Memory Hierarchy ---
\begin{frame}{Flash vs SRAM: The TinyML Memory Wall}
\begin{columns}[T]
\column{0.5\textwidth}
\textbf{Cloud GPU (H100)}
\begin{itemize}
  \item Weights in HBM (80\,GB)
  \item Activations in SRAM (50\,MB L2)
  \item Bandwidth: 3.35\,TB/s
  \item Model $\leq$ 80\,GB $\Rightarrow$ fits
\end{itemize}

\column{0.5\textwidth}
\textbf{TinyML MCU (ESP32-S3)}
\begin{itemize}
  \item Weights in \textbf{Flash} (8\,MB)
  \item Activations in \textbf{SRAM} (512\,KiB)
  \item Flash BW: 80\,MB/s
  \item Model $\leq$ 8\,MB $\Rightarrow$ fits
  \item Model $\leq$ 512\,KiB $\Rightarrow$ SRAM-only (12$\times$ faster)
\end{itemize}
\end{columns}

\vspace{1em}
\textbf{Key insight:} mlsysim automatically selects the right memory tier:
\begin{enumerate}
  \item Model fits in SRAM $\Rightarrow$ use SRAM bandwidth
  \item Model in Flash $\Rightarrow$ use Flash bandwidth (bottleneck!)
  \item Model exceeds Flash $\Rightarrow$ infeasible
\end{enumerate}
\end{frame}

% --- 7.4 Energy per Inference ---
\begin{frame}{Energy per Inference: $\mu$J to Joules}
\centering
\begin{tabular}{lrrl}
\toprule
\textbf{Device} & \textbf{TDP} & \textbf{Energy/Inf} & \textbf{Use Case} \\
\midrule
nRF52840       & 15\,mW  & $\sim$50\,$\mu$J & Keyword spotting \\
ESP32-S3       & 400\,mW & $\sim$1\,mJ      & Person detection \\
Jetson Orin NX & 25\,W   & $\sim$25\,mJ     & Object detection \\
\rowcolor{gray!15}
H100 SXM       & 700\,W  & $\sim$50\,J      & LLM inference \\
\bottomrule
\end{tabular}

\vspace{1em}
\begin{itemize}
  \item 6 orders of magnitude in energy per inference
  \item TinyML: battery life is the primary constraint
  \item Duty cycling (sleep 99\% of the time) extends battery to years
\end{itemize}
\end{frame}

% --- 7.5 Live Demo: Hardware Comparison ---
\begin{frame}[fragile]{Live Demo: nRF52840 vs ESP32 vs H100}
\begin{minted}[fontsize=\scriptsize]{python}
import mlsysim

# A tiny keyword-spotting model (~250 KB)
model = mlsysim.Models.Tiny.KeywordSpotting

devices = [
    mlsysim.Hardware.Tiny.nRF52840,
    mlsysim.Hardware.Tiny.ESP32_S3,
    mlsysim.Hardware.Cloud.H100,
]

for hw in devices:
    p = mlsysim.Engine.solve(model, hw, batch_size=1,
                             precision="int8", efficiency=0.3)
    energy_str = f"{p.energy.to('uJ'):.1f~P}" if p.energy.magnitude < 0.001 \
                 else f"{p.energy.to('mJ'):.2f~P}"
    print(f"{hw.name:>30}: {p.latency.to('ms'):>10.2f~P} | "
          f"{p.bottleneck:<8} | Energy: {energy_str}")
\end{minted}

\vspace{0.5em}
\begin{exampleblock}{Observation}
nRF52840 is memory-bound (Flash bottleneck).\\
H100 finishes in microseconds---but wastes 700\,W doing it.\\
\textbf{Right-sizing hardware matters.}
\end{exampleblock}
\end{frame}

% --- 7.6 Same Roofline, Different Physics ---
\begin{frame}{Same Roofline, Different Physics}
\begin{center}
\large
\[
\text{Latency} = \max\!\left(
  \frac{\text{FLOPs}}{\text{Peak} \times \eta},\;
  \frac{\text{Weights}}{\text{BW}}
\right) + \text{Overhead}
\]
\end{center}

\vspace{0.5em}
\begin{tabular}{lll}
\toprule
\textbf{Term} & \textbf{Cloud} & \textbf{TinyML} \\
\midrule
Peak        & 989 TFLOPS (tensor cores) & 500 MFLOPS (scalar ALU) \\
BW          & 3.35 TB/s (HBM3)         & 80 MB/s (SPI Flash) \\
Overhead    & 0.01\,ms (CUDA dispatch)  & 1.0\,ms (TFLite Micro) \\
Bottleneck  & Memory (LLM inference)    & Memory (Flash read) \\
\bottomrule
\end{tabular}

\vspace{0.5em}
\textbf{Both are memory-bound at batch size 1.}\\
The same equation diagnoses the same bottleneck.
\end{frame}

% --- 7.7 Exercise ---
\begin{frame}[fragile]{Exercise: Right-Sizing for IoT}
\begin{alertblock}{Your Task (5 minutes)}
A smart doorbell needs person detection at $<$100\,ms latency on a
coin-cell battery (500\,mAh @ 3V = 1.5\,Wh). Compare ESP32-S3 vs
nRF52840: which device can run inference for a full year at 1 inference/minute?
\end{alertblock}

\begin{minted}[fontsize=\scriptsize]{python}
import mlsysim

model = mlsysim.Models.Tiny.PersonDetection
for hw in [mlsysim.Hardware.Tiny.ESP32_S3,
           mlsysim.Hardware.Tiny.nRF52840]:
    p = mlsysim.Engine.solve(model, hw, precision="int8", efficiency=0.3)
    energy_j = p.energy.to("J").magnitude
    inferences_per_year = 60 * 24 * 365
    # YOUR CODE: compute total energy, compare to 1.5 Wh battery
\end{minted}

\note{Expected answer: nRF52840 at 15\,mW uses far less energy per inference.
With duty cycling (awake only during inference), nRF52840 lasts $>$1 year.
ESP32-S3 at 400\,mW drains the battery in weeks without aggressive duty cycling.}
\end{frame}

% --- 7.8 Takeaway ---
\begin{frame}{Key Takeaway: TinyML to Frontier}
\begin{center}
\Large
\begin{enumerate}
  \item The Roofline model spans \textbf{9 orders of magnitude}
  \item \textbf{Right-size} hardware to the workload, not the other way around
  \item At TinyML scale, \textbf{energy and memory} dominate; at cloud scale, \textbf{cost and communication}
\end{enumerate}

\vspace{1em}
\normalsize
\textit{``The best accelerator is the smallest one\\
that meets your latency and accuracy SLA.''}
\end{center}
\end{frame}


% =====================================================================
% PART 8: ADVANCED TOPICS
% =====================================================================
\section{Advanced Topics}

% --- 8.1 Key Question ---
\begin{frame}{Key Question}
\begin{center}
\Large\bfseries
How do you compose multiple analytical models\\[4pt]
into an end-to-end system analysis?
\end{center}
\vfill
\small
Single-wall analysis answers ``what is the bottleneck?''\\
\textbf{Pipeline composition} answers ``what is the total cost of the whole stack?''
\end{frame}

% --- 8.2 Pipeline Composition Pattern ---
\begin{frame}[fragile]{The Pipeline Composition Pattern}
\begin{minted}[fontsize=\scriptsize]{python}
from mlsysim.core.pipeline import Pipeline
from mlsysim.core.solver import (
    ScalingModel, DistributedModel, EconomicsModel)

# Chain three resolvers: Algorithm -> Fleet -> Economics
pipe = Pipeline([
    ScalingModel(),        # Wall 11: How much compute?
    DistributedModel(),    # Wall 14: Communication overhead
    EconomicsModel(),      # Wall 17: Total cost
])

# Inspect the DAG before running
print(pipe.explain())
\end{minted}

\vspace{0.5em}
\texttt{explain()} shows:
\begin{itemize}
  \item What each stage requires and produces
  \item Which walls from the 22-wall taxonomy are covered
  \item Where gaps exist (missing inputs)
\end{itemize}
\end{frame}

% --- 8.3 Pipeline Run ---
\begin{frame}[fragile]{Running the Pipeline}
\begin{minted}[fontsize=\scriptsize]{python}
from mlsysim.core.constants import Q_

results = pipe.run(
    compute_budget=Q_("1e21 FLOP"),
    fleet=mlsysim.Systems.Clusters.Research_256,
    duration_days=30,
    datacenter=mlsysim.Infra.Grids.Quebec,
    mfu=0.45,
    infrastructure_multiplier=2.5
)

# Results keyed by resolver class name
for stage_name, result in results.items():
    print(f"\n--- {stage_name} ---")
    print(result)
\end{minted}

\vspace{0.5em}
\textbf{Key design principle:} Each resolver's output fields become available
as inputs to subsequent stages. The pipeline is \emph{not} a black box---it
is a transparent analytical tool.
\end{frame}

% --- 8.4 Sensitivity Analysis ---
\begin{frame}[fragile]{Wall 21: Sensitivity Analysis}
\begin{minted}[fontsize=\scriptsize]{python}
import mlsysim

solver = mlsysim.SensitivitySolver()
result = solver.solve(
    model=mlsysim.Models.Language.Llama3_8B,
    hardware=mlsysim.Hardware.H100,
    precision="fp16",
    perturbation_pct=10.0,  # +10% perturbation
    efficiency=0.5
)

print(f"Binding constraint: {result.binding_constraint}")
print(f"Baseline latency:   {result.baseline_latency:~P}")
print(f"\nSensitivities (dT/T per +10%):")
for param, sensitivity in result.sensitivities.items():
    arrow = "<<<" if param == result.binding_constraint else ""
    print(f"  {param:>20}: {sensitivity:+.4f}  {arrow}")
\end{minted}

\vspace{0.3em}
\small
\textbf{Rule:} Invest in the parameter with the \emph{largest} $|\partial T / \partial p|$.\\
Improving a non-binding parameter yields \textbf{zero} measurable gain.
\end{frame}

% --- 8.5 Inverse Roofline (Synthesis) ---
\begin{frame}[fragile]{Wall 22: Inverse Roofline (SynthesisSolver)}
\begin{minted}[fontsize=\scriptsize]{python}
import mlsysim
from mlsysim.core.constants import Q_

solver = mlsysim.SynthesisSolver()
result = solver.solve(
    model=mlsysim.Models.Language.Llama3_8B,
    target_latency=Q_("50 ms"),
    precision="fp16",
    efficiency=0.5
)

print(f"To serve Llama-3-8B in <50ms:")
print(f"  Required BW:     {result.required_bw:.0f~P}")
print(f"  Required FLOPS:  {result.required_flops:.1f~P}")
print(f"  Required Memory: {result.required_memory:.1f~P}")
\end{minted}

\vspace{0.5em}
\textbf{Usage:} ``I need 50\,ms TTFT. What hardware do I need?''\\
$\Rightarrow$ Derive specs from SLA, then match to real accelerators.
\end{frame}

% --- 8.6 Fallacies & Pitfalls ---
\begin{frame}{Fallacies \& Pitfalls (Patterson Tradition)}
\begin{block}{Fallacy: ``Cheaper hardware is always more cost-effective''}
Reality: A 2$\times$ cheaper GPU that takes 3$\times$ longer has \emph{higher} TCO.\\
$\text{TCO} = \text{CapEx} + \text{OpEx}$; slower hardware burns more electricity.
\end{block}

\begin{block}{Fallacy: ``More FLOPS = proportionally faster''}
Reality: A100 $\to$ H100 is 6.3$\times$ peak FLOPS but only $\sim$2$\times$ for
memory-bound LLM inference. The binding constraint is bandwidth, not compute.
\end{block}

\begin{block}{Pitfall: Optimizing a non-binding parameter}
If inference is memory-bound, doubling FLOPS gives 0\% speedup.\\
Use \texttt{SensitivitySolver} to find the binding constraint \emph{first}.
\end{block}

\begin{block}{Pitfall: Ignoring the infrastructure multiplier}
GPU cost is 40\% of total datacenter TCO. Budgeting for GPUs alone\\
underestimates the true cost by 2--3$\times$.
\end{block}
\end{frame}

% --- 8.7 The Iron Law Revisited ---
\begin{frame}{The Iron Law Revisited: All Five Terms}
\begin{center}
\Large
\[
T = \frac{\text{FLOPs}}{N \times \text{Peak} \times \text{MFU} \times \eta_{\text{scale}} \times \text{Goodput}}
\]
\end{center}

\vspace{0.5em}
\centering
\begin{tabular}{lll}
\toprule
\textbf{Term} & \textbf{What reduces it} & \textbf{Wall} \\
\midrule
$N$                  & Budget (buy more GPUs)        & Wall 17 \\
Peak                 & GPU generation (H100$\to$B200)& Wall 1 \\
MFU                  & FlashAttention, kernel fusion  & Wall 3 \\
$\eta_{\text{scale}}$& Network BW, gradient compression & Wall 14 \\
Goodput              & Checkpointing, fault tolerance & Walls 15, 19 \\
\bottomrule
\end{tabular}

\vspace{0.5em}
\textbf{Every wall in the 22-wall taxonomy maps to one of these five terms.}\\
That is the entire framework.
\end{frame}

% --- 8.8 Exercise ---
\begin{frame}[fragile]{Exercise: Binding Constraint Analysis}
\begin{alertblock}{Your Task (5 minutes)}
Use \texttt{SensitivitySolver} on Llama-3-8B + H100.
Which parameter is binding? Now switch to \texttt{batch\_size=64} and
re-run. Does the binding constraint change?
\end{alertblock}

\begin{minted}[fontsize=\scriptsize]{python}
import mlsysim

solver = mlsysim.SensitivitySolver()

# Batch size 1
r1 = solver.solve(mlsysim.Models.Language.Llama3_8B,
                   mlsysim.Hardware.H100, efficiency=0.5)
print(f"bs=1:  binding = {r1.binding_constraint}")

# Batch size 64
p64 = mlsysim.Engine.solve(mlsysim.Models.Language.Llama3_8B,
                            mlsysim.Hardware.H100, batch_size=64)
print(f"bs=64: bottleneck = {p64.bottleneck}")
\end{minted}

\note{Expected answer: At batch\_size=1, memory\_bandwidth is binding
(memory-bound regime). At batch\_size=64, peak\_flops becomes binding
(compute-bound regime). This IS the Roofline transition.}
\end{frame}


% =====================================================================
% PART 9: WRAP-UP & FUTURE
% =====================================================================
\section{Wrap-up \& Future}

% --- 9.1 The 22 Walls in One Slide ---
\begin{frame}{The 22 Walls of ML Systems}
\centering
\small
\begin{tabular}{clcl}
\toprule
\textbf{\#} & \textbf{Wall} & \textbf{\#} & \textbf{Wall} \\
\midrule
1  & Compute         & 12 & Reasoning (Emerging) \\
2  & Memory          & 13 & Fidelity (Compression) \\
3  & Software (MFU)  & 14 & Communication \\
4  & Serving         & 15 & Fragility (Reliability) \\
5  & Batching (KV)   & 16 & Multi-tenant (Queueing) \\
6  & Streaming       & 17 & Capital (TCO) \\
7  & Tail Latency    & 18 & Sustainability \\
8  & Ingestion       & 19 & Checkpoint \\
9  & Transformation  & 20 & Safety (Privacy) \\
10 & Locality        & 21 & Sensitivity \\
11 & Complexity      & 22 & Synthesis (Inverse) \\
\bottomrule
\end{tabular}

\vspace{0.5em}
\textbf{6 Domains:} Node $\cdot$ Data $\cdot$ Algorithm $\cdot$ Fleet $\cdot$ Operations $\cdot$ Meta-Analysis

\vspace{0.3em}
\alert{Every wall maps to one term in the Iron Law.}
\end{frame}

% --- 9.2 What This Tool Does NOT Model ---
\begin{frame}{What \texttt{mlsysim} Does \emph{Not} Model}
\begin{columns}[T]
\column{0.5\textwidth}
\textbf{Not modeled (v0.1.0):}
\begin{itemize}
  \item Cache effects / tiling strategies
  \item Operator fusion decisions
  \item Real network congestion
  \item CUDA kernel scheduling
  \item Non-deterministic jitter
  \item Multi-model co-location
  \item Memory fragmentation
  \item Compiler optimizations
\end{itemize}

\column{0.5\textwidth}
\textbf{By design:}
\begin{itemize}
  \item Not a cycle-accurate simulator
  \item Not a profiler replacement
  \item Not a deployment tool
  \item \textbf{Is:} a first-principles\\
        analytical reasoning tool
\end{itemize}

\vspace{0.5em}
\textbf{Use it for:}
\begin{itemize}
  \item ``Which constraint is binding?''
  \item ``Is this hardware sufficient?''
  \item ``What should I try first?''
\end{itemize}
\end{columns}

\vspace{0.5em}
\centering
\small\itshape
``All models are wrong, but some are useful.'' --- George Box
\end{frame}

% --- 9.3 v0.2.0 Roadmap ---
\begin{frame}{v0.2.0 Roadmap}
\begin{enumerate}
  \item \textbf{MoE support} --- Mixture-of-Experts routing, expert parallelism
  \item \textbf{Multi-query / Grouped-query attention} --- GQA KV cache modeling
  \item \textbf{Network congestion model} --- Contention-aware collective timing
  \item \textbf{Spot instance pricing} --- Cloud cost optimization with preemption
  \item \textbf{YAML/JSON scenario files} --- Declarative workload specification
  \item \textbf{Interactive Marimo dashboards} --- Live parameter exploration
  \item \textbf{HuggingFace Hub integration} --- Auto-import model architectures
  \item \textbf{Calibration suite} --- Validate against published MLPerf results
\end{enumerate}

\vspace{0.5em}
\begin{center}
Contributions welcome: \texttt{github.com/harvard-edge/mlsysim}
\end{center}
\end{frame}

% --- 9.4 Design Challenge (Capstone) ---
\begin{frame}[fragile]{Design Challenge: Capstone Exercise}
\begin{alertblock}{The Problem}
\textbf{\$5M budget.} Serve Llama-3 70B at \textbf{1{,}000 QPS} with
\textbf{$<$100\,ms TTFT} in \textbf{two regions} (US-East + EU-West).
Design the fleet.
\end{alertblock}

\vspace{0.3em}
\textbf{You must specify:}
\begin{enumerate}
  \item Hardware choice (which GPU? how many?)
  \item Parallelism strategy (TP $\times$ PP per replica)
  \item Precision (FP16? INT8? FP8?)
  \item Geographic placement (which grids? carbon impact?)
  \item Redundancy (how many replicas per region?)
\end{enumerate}

\vspace{0.3em}
\textbf{Tools at your disposal:}
\begin{itemize}
  \item \texttt{Engine.solve()} --- single-node feasibility
  \item \texttt{EconomicsModel.solve()} --- TCO
  \item \texttt{SustainabilityModel.solve()} --- carbon
  \item \texttt{SensitivitySolver.solve()} --- binding constraint
  \item \texttt{SynthesisSolver.solve()} --- required specs from SLA
\end{itemize}

\note{This is intentionally under-specified. Students must make and defend
assumptions. Key insight: 70B at FP16 = 140\,GB, needs tensor parallelism
across 2+ H100s per replica. At 1000 QPS with 100\,ms TTFT, you need many
replicas. Budget constrains how many.}
\end{frame}

% --- 9.5 Resources & Q&A ---
\begin{frame}{Resources \& Next Steps}
\begin{columns}[T]
\column{0.55\textwidth}
\textbf{Get Started}
\begin{itemize}
  \item \texttt{pip install mlsysim}
  \item GitHub: \texttt{harvard-edge/mlsysim}
  \item 14 examples + 3 Marimo notebooks
  \item Full docs: \texttt{mlsysim.readthedocs.io}
\end{itemize}

\vspace{0.5em}
\textbf{The Textbook}
\begin{itemize}
  \item \emph{Machine Learning Systems}
  \item Volume I: Foundations (single node)
  \item Volume II: Systems at Scale (fleet)
  \item \texttt{mlsysbook.ai}
\end{itemize}

\column{0.40\textwidth}
\textbf{Key Papers}
\begin{itemize}
  \item Williams et al.\ (2009)\\
        {\footnotesize Roofline Model}
  \item Chowdhery et al.\ (2022)\\
        {\footnotesize PaLM / MFU}
  \item Hoffmann et al.\ (2022)\\
        {\footnotesize Chinchilla Scaling}
  \item Patterson et al.\ (2021)\\
        {\footnotesize Carbon \& Training}
  \item Kwon et al.\ (2023)\\
        {\footnotesize PagedAttention}
\end{itemize}
\end{columns}

\vspace{1em}
\begin{center}
\Large\bfseries Thank you! Questions?
\end{center}
\end{frame}
